\begin{proof}
We first parametrize the entire feasible set, then remove the phase
symmetries at its limiting maximizing set.  The resulting maximum is a
nondegenerate analytic critical point in seven real variables.

\emph{An exact compact parametrization.}
Put $t=s^2$, initially allowing either sign of the real parameter $s$.
Write the first column of $V$ as $u=sz$, where
$z\in\FC^3$ and $\Fnorm{z}=1$, and write the second column as $v$.
For $|s|<1$, the Schur complement of $I-V^*V$ gives the exact equivalence
\begin{equation}
 V^*V\preceq I
 \quad\Longleftrightarrow\quad
 \Fnorm{v}^2+\frac{s^2}{1-s^2}|z^*v|^2\le1.
 \label{I:F:eq:f2-column-ellipsoid}
\end{equation}
Thus every feasible second column is uniquely of the form
\begin{equation}
 v=W_s(z)\nu,\qquad \Fnorm{\nu}\le1,\qquad
 W_s(z)=I_3-\bigl(1-\sqrt{1-s^2}\bigr)zz^*.
 \label{I:F:eq:f2-W-parametrization}
\end{equation}
The matrix $W_s(z)$ is analytic and even in $s$ near zero.  In these
coordinates,
\[
 B=\begin{pmatrix}sz_1&v_1\\sz_2&v_2\end{pmatrix},\qquad
 \ell=(sz_3,v_3),
\]
and direct multiplication gives
\begin{align}
 \ell B^2e_1&=s\,\mathcal J(s,z,\nu),\notag\\
 \mathcal J(s,z,\nu)
 &=z_2v_2v_3+s z_2(z_3v_1+z_1v_3)+s^2z_3z_1^2.
 \label{I:F:eq:f2-J-objective}
\end{align}
Consequently
\begin{equation}
 h(s^2)=s^2\mathcal M(s)^2,\qquad
 \mathcal M(s)=\max_{\Fnorm{z}=1,\,\Fnorm{\nu}\le1}
                  |\mathcal J(s,z,\nu)|.
 \label{I:F:eq:f2-M-definition}
\end{equation}

At $s=0$,
\[
 |\mathcal J(0,z,\nu)|=|z_2\nu_2\nu_3|\le\frac12.
\]
Equality is possible exactly on the phase torus
\begin{equation}
 z_1=z_3=\nu_1=0,\qquad
 |z_2|=1,\qquad |\nu_2|=|\nu_3|=2^{-1/2}.
 \label{I:F:eq:f2-maximizing-torus}
\end{equation}
The parameter space is compact, and the objective converges uniformly
as $s\to0$.  Hence every global maximizing pair for sufficiently small
$|s|$ lies arbitrarily close to this torus.

\emph{All complex phase degeneracies can be removed.}
The original optimization is invariant under
\begin{equation}
 B\longmapsto e^{i\theta}U^*BU,\qquad
 \ell\longmapsto e^{i\phi}\ell U,\qquad
 U=\operatorname{diag}(1,e^{i\psi}).
 \label{I:F:eq:f2-phase-action}
\end{equation}
The scalar $\ell B^2e_1$ is multiplied only by
$e^{i(\phi+2\theta)}$.  On the compact parameters, the same action is
\[
 z\longmapsto Dz,\qquad
 \nu\longmapsto e^{i\psi}D\nu,\qquad
 D=\operatorname{diag}\bigl(e^{i\theta},e^{i(\theta-\psi)},e^{i\phi}\bigr).
\]
Indeed \eqref{I:F:eq:f2-W-parametrization} is equivariant under diagonal
unitaries.  The phases of $(z_2,\nu_2,\nu_3)$ change respectively by
$(\theta-\psi,\theta,\phi+\psi)$, an invertible linear transformation
of the three phase parameters.  Since those three coordinates are
nonzero near \eqref{I:F:eq:f2-maximizing-torus}, they may all be fixed to be
positive real in a smooth local gauge.  No transverse complex direction
is excluded by this choice.

A global maximizer also satisfies $\Fnorm{\nu}=1$ for all sufficiently
small $|s|$.  To see this, use radial coordinates
$\nu=\rho\widehat\nu$, $\Fnorm{\widehat\nu}=1$.
At the limiting torus the derivative of
$|\mathcal J|$ in the increasing radial direction is $1$ at $\rho=1$.
It remains positive in a neighborhood of that compact torus.  Global
maximizers approach the torus, so a maximizing point with $\rho<1$
would admit a feasible radial increase of its objective.

\emph{Seven real coordinates and a strictly negative Hessian.}
After fixing the phases and the radial constraint, choose local
coordinates
\begin{equation}
 \begin{gathered}
 z_1,z_3,\nu_1\in\FC,\qquad \xi\in\FR,\qquad
 z_2=\sqrt{1-|z_1|^2-|z_3|^2},\\
 \nu_2=\sqrt{1-|\nu_1|^2}\cos(\pi/4+\xi),\qquad
 \nu_3=\sqrt{1-|\nu_1|^2}\sin(\pi/4+\xi).
 \end{gathered}
 \label{I:F:eq:f2-seven-coordinates}
\end{equation}
Let $x\in\FR^7$ list, in order,
$\Re z_1,\Im z_1,\Re z_3,\Im z_3,\Re\nu_1,\Im\nu_1,\xi$.
The gauge sends the entire limiting maximizing torus to $x=0$.
The objective
$F(s,x)=|\mathcal J(s,z(x),\nu(x))|$ is real analytic near $(0,0)$,
because $\mathcal J(0,0)=1/2\ne0$.
At $s=0$ its expansion is
\begin{equation}
 F(0,x)=\frac12-\frac{|z_1|^2+|z_3|^2}{4}
                  -\frac{|\nu_1|^2}{2}-\xi^2+O(\Fnorm{x}^4).
 \label{I:F:eq:f2-transverse-Hessian}
\end{equation}
Thus
\begin{equation}
 \partial_x^2F(0,0)=
 \operatorname{diag}\!\left(-\frac12,-\frac12,-\frac12,-\frac12,
                           -1,-1,-2\right)\prec0.
 \label{I:F:eq:f2-Hessian-matrix}
\end{equation}
The full expansion through total degree two is
\begin{align}
 F(s,x)=\frac12-\frac14s^2+\frac{s\,\Re z_1}{\sqrt2}
          -\frac{|z_1|^2+|z_3|^2}{4}
          -\frac{|\nu_1|^2}{2}-\xi^2
          +O\bigl((|s|+\Fnorm{x})^3\bigr).
 \label{I:F:eq:f2-full-quadratic}
\end{align}
For completeness, the pure $s^2$ term follows from the exact identity
$\mathcal J(s,0)=\sqrt{1-s^2}/2$.  In
\eqref{I:F:eq:f2-J-objective}, the only mixed term of total degree two
contributing to the modulus is $s\Re z_1/\sqrt2$.

The analytic implicit function theorem applied to
$\partial_xF(s,x)=0$ gives a unique analytic critical branch $x=x(s)$
near zero.  The negative definite Hessian persists in a fixed small
coordinate neighborhood, where the objective is strictly concave.
Every global maximizer lies in that neighborhood after the phase gauge,
by the compact localization already proved.  It must therefore be this
critical branch.  This justifies applying the local analytic expansion
to the \emph{global} maximum $\mathcal M(s)$.
Differentiating the critical equations using
\eqref{I:F:eq:f2-Hessian-matrix}--\eqref{I:F:eq:f2-full-quadratic} gives
\begin{equation}
 x'(0)=(\sqrt2,0,0,0,0,0,0).
 \label{I:F:eq:f2-critical-derivative}
\end{equation}
Substituting into \eqref{I:F:eq:f2-full-quadratic} yields
\begin{equation}
 \mathcal M(s)=\frac12+\frac14s^2+O(s^3).
 \label{I:F:eq:f2-M-preliminary}
\end{equation}

\emph{Physical evenness improves the remainder.}
The feasible physical matrices with first-column squared norm $s^2$
are identical for $s$ and $-s$.  Equivalently, replacing
$(s,z,\nu)$ by $(-s,-z,\nu)$ preserves $v$ and changes
$\mathcal J$ only by its sign.  Therefore
$\mathcal M(-s)=\mathcal M(s)$.  We have proved that the maximizing
value is analytic near zero, so all its odd Taylor coefficients vanish.
Equation \eqref{I:F:eq:f2-M-preliminary} sharpens to
\begin{equation}
 \mathcal M(s)=\frac12+\frac14s^2+O(s^4).
 \label{I:F:eq:f2-M-even}
\end{equation}
Finally \eqref{I:F:eq:f2-M-definition} gives
$h(t)=t(1/2+t/4+O(t^2))^2$, which is
\eqref{I:F:eq:f2-h-expansion}.
\end{proof}
